\appendix
\section{Corpus composition}
\label{app:corpus}

Table~\ref{tab:corpus} summarises the corpus structure and lists one sample
sentence per stratum. The four domains are constructed to differ in register,
vocabulary, and numerical density. The broadcast-news domain uses formal
written register with institutional vocabulary, the conversation domain uses
colloquial everyday phrasing, the technical domain uses scientific and
engineering terminology, and the numerals-and-entities domain concentrates
dates, times, quantities, prices, addresses, codes, and named entities, with
all numbers written out in words as described in
Section~\ref{sec:methods}. Length strata are defined per language as short,
medium, and long as given in Section~\ref{sec:methods}, and every
domain-by-length cell contains twenty sentences per language, giving 240
sentences per language and 480 in total. The severity sweep subset of 144
sentences is sampled evenly across all cells, six sentences per cell per
language.

\begin{table}[h]
\centering
\caption{Corpus structure with one sample sentence per domain.}
\label{tab:corpus}
\small
\begin{tabular}{lp{9.3cm}}
\toprule
Domain & Sample sentence \\
\midrule
News (zh) & 国家统计局今天发布了上月经济运行数据。 \\
News (en) & The central bank raised interest rates by twenty five basis points to curb inflation. \\
Conversation (zh) & 我打算下班之后去健身房锻炼一个小时。 \\
Conversation (en) & I was thinking we could drive to the beach this weekend if the weather holds. \\
Technical (zh) & 混响是声波在房间内多次反射形成的，会显著降低远场语音识别的准确率。 \\
Technical (en) & Reverberation arises when sound waves reflect repeatedly off room surfaces. \\
Numerals (zh) & 从北京到上海的高铁大约需要四个半小时。 \\
Numerals (en) & The high speed train from Beijing to Shanghai takes about four and a half hours. \\
\bottomrule
\end{tabular}
\end{table}

\section{Qualitative error analysis}
\label{app:qualitative}

Table~\ref{tab:examples} collects representative model responses that
illustrate the failure patterns behind the aggregate numbers in
Section~\ref{sec:results}.

The refusal behaviour of Section~\ref{sec:results:main} is shown in the
first row: under 0~dB babble the Qwen-Audio recogniser returns an empty
transcript for a substantial English sentence, which the benchmark scores as
a complete failure. The second row shows the opposite regime: the same
system, on a long English news sentence at the same severity, produces a
verbatim transcript, which underlines how bimodal the sentence-level
distributions of Fig.~\ref{fig:dist} are. Row three shows the dominant error
mode of the specialised streaming recogniser under severe noise, where most
of the sentence is lost and only a short fragment survives. Row four shows a
typical near-miss of the LLM recogniser in Chinese, where a single
content word, density, is replaced by another word of similar phonetic
shape, speed, without any structural damage. Row five shows the collapse of
the telephone-band specialist under reverberation, which occasionally
produces short disconnected fragments rather than transcripts. Row six
shows an error that the recognition stage never sees but the translation
consumer experiences: under 0~dB noise the omni end-to-end path substitutes
a plausible but wrong entity, rendering the sports centre of the source as a
grand theatre, a hallucination that a string-matching metric penalises only
lightly but that is catastrophic for factual fidelity. Row seven shows the
narrowband immunity of the LLM recogniser, which returns a verbatim
transcript of telephone-band audio. Together the examples support the
reading of Section~\ref{sec:discussion}: the systems differ less in whether
they fail than in how, and the failure modes, from verbatim robustness to
silent entity substitution to outright refusal, have distinct downstream
consequences.

\begin{table}[h]
\centering
\caption{Representative system responses under degradation.}
\label{tab:examples}
\small
\begin{tabular}{llp{6.1cm}}
\toprule
Condition & System & Reference and response \\
\midrule
0 dB, en & Qwen-Audio-3.1 & Reference: Police arrested two suspects on Tuesday. Response: empty transcript. \\
0 dB, en & Qwen-Audio-3.1 & Reference: The transport minister told reporters on Friday that the government would invest heavily in electric vehicle charging infrastructure. Response: verbatim transcript. \\
0 dB, en & Paraformer-RT-v2 & Reference: The battery lasts twenty hours. Response: About. \\
0 dB, zh & Qwen-Audio-3.1 & Reference: 这款处理器的晶体管密度提升了近一倍。 Response: 这款处理器的晶体管速度提升了近一倍。 \\
Reverb, zh & Paraformer-RT-8K & Reference: 记者从发布会上了解到，本次合作项目将覆盖全国二十多个主要城市。 Response: 七月同。非常了解到。 \\
0 dB, E2E & Omni-Flash & Source: 新建成的体育中心可容纳六万名观众。 Translation: The newly built Grand Theatre can accommodate 60,000 audience members. \\
Narrow, en & Qwen-Audio-3.1 & Reference: This coffee shop is cozy. Response: verbatim transcript. \\
$-10$ dB, en & Qwen-Audio-3.1 & Reference: The battery lasts twenty hours. Response: fluent unrelated Chinese text of several sentences (hallucination). \\
\bottomrule
\end{tabular}
\end{table}

\section{Cross-system text normalisation}
\label{app:normalisation}

Comparing transcripts across commercial recognisers requires a normalisation
protocol that neutralises provider-specific inverse text normalisation. The
protocol applies the following rules identically to references and
hypotheses, in order. Unicode compatibility folding and case folding are
applied first, and punctuation is removed. Numeric expressions are then
canonicalised to digit form on both sides. For Chinese, the rules handle
cardinal numbers with the standard multipliers, digit-by-digit readings such
as room and code numbers where each digit is read separately, decimal
fractions written with the point particle, time-of-day expressions in which
the hour and minute parts are converted separately, percentage markers,
which are deleted on both sides, and scale multipliers, where a digit
followed by wan, yi, or their combinations is expanded to the full
multiplied-out form, so that twenty-six million written as 2600万 in one
system's output and as 26000000 in another's both canonicalise identically.
For English, the rules handle cardinal numbers including compound
hyphenated forms and the and connector, time and digit-by-digit readings
such as twelve zero eight and nine forty, which canonicalise to 1208 and
940, decimal fractions expressed with point, scale multipliers from thousand
to trillion, and ordinal expressions of both word and digit form, which are
masked to a common placeholder since their spoken-to-written realisations
differ across systems in ways that carry no acoustic information. The rule
set was developed and frozen by inspecting all high-error sentences of the
clean condition until no residual mismatches attributable to normalisation
remained, a procedure that could in principle favour whichever system's
conventions dominate the clean outputs, so a sensitivity analysis was run:
re-scoring the main grid with case and punctuation folding only, dropping
numeral canonicalisation entirely, raises clean-condition error rates from
0.5, 2.3, and 0.3 percent to 6.1, 6.2, and 0.8 percent for Qwen-Audio,
Paraformer, and Qwen3-ASR respectively, while the 0~dB scores move by less
than 1.5 points and the system ordering is unchanged in every condition.
The normalisation therefore compresses but does not create the reported
system differences, and the residual protocol noise on clean speech is
bounded above by the strongest system's per-language clean rate of 0.14
percent in Mandarin (0.29 percent pooled over both languages), with the
caveat that the two systems most affected by numeral conventions differ
most under the raw protocol. The full normalisation code is released with
the benchmark.

\section{Complete pairwise significance results}
\label{app:significance}

Table~\ref{tab:pairwise} reports the pairwise Wilcoxon signed-rank
comparisons between recognition systems that remain significant at the
Holm-corrected level within the family of all pairwise comparisons of the
same condition and language, so the table lists only significant results
by construction. Two protocol notes are needed to read it correctly.
First, comparisons among the three general-purpose recognisers are computed
on the full 480-sentence main grid, 240 sentences per language, whereas
every comparison involving the omni system is computed on the 144-sentence
sweep subset, 72 per language, because that system never ran the full grid;
the two blocks of the table are same-subset within themselves, and the
omni rows state their subset explicitly. Second, the fraction of zero
paired differences is substantial in mild conditions, 72 to 83 percent
under reverberation and narrowband, so significance there is carried by
the minority of sentences where the systems differ. The consistency of
the ordering is itself informative: within every condition and language,
the two LLM-ASR systems occupy the leading positions and the specialised
streaming recogniser the trailing position, with the exception that
Qwen3-ASR leads Qwen-Audio 3.1 in English under reverberation, clipping,
narrowband, and 0~dB noise, where the refusal behaviour of the latter
penalises it, while the omni system is statistically inseparable from the
Qwen systems in English under noise and reverberation on its subset.

\begin{table}[h]
\centering
\caption{Significant Holm-corrected pairwise Wilcoxon comparisons by
condition and language. Cells give the mean error rates in percent of the
two systems compared. Upper block: full main grid, $n{=}240$ per language.
Lower block: 144-sentence sweep subset, $n{=}72$ per language (omni
system).}\label{tab:pairwise}
\small
\setlength{\tabcolsep}{4pt}
\begin{tabular}{llll}
\toprule
Condition, language & Comparison & Means (\%) & $p$ \\
\midrule
0 dB, zh & Qwen-Audio vs Paraformer-RT & 6.3 vs 41.6 & $<0.001$ \\
0 dB, zh & Qwen3-ASR vs Paraformer-RT & 13.1 vs 41.6 & $<0.001$ \\
0 dB, zh & Qwen-Audio vs Qwen3-ASR & 6.3 vs 13.1 & $<0.001$ \\
0 dB, en & Qwen-Audio vs Paraformer-RT & 50.2 vs 75.5 & $<0.001$ \\
0 dB, en & Qwen3-ASR vs Paraformer-RT & 37.1 vs 75.5 & $<0.001$ \\
0 dB, en & Qwen-Audio vs Qwen3-ASR & 50.2 vs 37.1 & $<0.001$ \\
Reverb, zh & Qwen-Audio vs Paraformer-RT & 0.2 vs 1.3 & $<0.001$ \\
Reverb, zh & Qwen3-ASR vs Paraformer-RT & 0.1 vs 1.3 & $<0.001$ \\
Reverb, en & Qwen3-ASR vs Paraformer-RT & 0.5 vs 5.1 & $<0.001$ \\
Reverb, en & Qwen-Audio vs Qwen3-ASR & 1.2 vs 0.5 & $0.011$ \\
Narrow, zh & Qwen-Audio vs Paraformer-RT & 0.1 vs 0.7 & $<0.001$ \\
Narrow, zh & Qwen3-ASR vs Paraformer-RT & 0.1 vs 0.7 & $<0.001$ \\
Narrow, en & Qwen3-ASR vs Paraformer-RT & 0.5 vs 3.6 & $<0.001$ \\
Narrow, en & Qwen-Audio vs Paraformer-RT & 1.0 vs 3.6 & $<0.001$ \\
Narrow, en & Qwen-Audio vs Qwen3-ASR & 1.0 vs 0.5 & $0.015$ \\
Clip, en & Qwen3-ASR vs Paraformer-RT & 0.4 vs 4.4 & $<0.001$ \\
Clip, en & Qwen-Audio vs Paraformer-RT & 1.6 vs 4.4 & $<0.001$ \\
Clip, en & Qwen-Audio vs Qwen3-ASR & 1.6 vs 0.4 & $<0.001$ \\
\midrule
\multicolumn{4}{l}{\footnotesize Sweep-subset comparisons involving the omni system ($n{=}72$ per language):} \\
0 dB, zh & Qwen-Audio vs Omni & 8.5 vs 20.2 & $<0.001$ \\
0 dB, zh & Qwen3-ASR vs Omni & 11.2 vs 20.2 & $<0.001$ \\
0 dB, zh & Paraformer-RT vs Omni & 42.4 vs 20.2 & $<0.001$ \\
0 dB, en & Qwen3-ASR vs Omni & 34.1 vs 43.3 & $0.018$ \\
0 dB, en & Paraformer-RT vs Omni & 76.1 vs 43.3 & $<0.001$ \\
Reverb, zh & Paraformer-RT vs Omni & 1.3 vs 0.1 & $0.003$ \\
Reverb, en & Paraformer-RT vs Omni & 5.0 vs 2.2 & $0.002$ \\
\bottomrule
\end{tabular}
\end{table}

\section{Numeric results of the severity sweep}
\label{app:sweep}

Table~\ref{tab:sweep} reports the complete error rates of the eight-level
severity sweep with confidence intervals, and Table~\ref{tab:ltfrontend}
reports the internal frontend recogniser of the streaming translation
service on the same subset. Three details of the sweep data deserve note.
The confidence intervals at 0~dB and below are wide for every system,
spanning ten percentage points or more, which reflects the bimodal
sentence-level behaviour documented in Section~\ref{sec:results:latency}.
The crossover between the classes occurs between 5 and 0~dB for both
languages, and at 5~dB the specialised system already shows its first
significant departure from its clean score while the LLM-ASR systems
remain within 0.3 points of theirs in Mandarin. The internal frontend of
the streaming translation service has partial coverage in some sweep
conditions, 52 to 134 of 144 sessions, because its websocket transport
dropped sessions that did not recover within the retry budget; its zero
score at 5~dB rests on 52 completed sessions and is a small-sample
fluctuation rather than evidence of an operating point better than clean.
The sub-zero conditions, 2.5, $-5$, and $-10$~dB, were added in the
revision round with identically seeded babble realisations, and their
per-point sample sizes equal those of the original ladder, 72 sentences
per language per level.

\begin{table}[h]
\centering
\caption{Recognition error rates in percent on the severity sweep subset,
144 sentences per level, with bootstrap 95 percent confidence intervals.
Values above 100 percent arise from insertion-dominated outputs.}
\label{tab:sweep}
\footnotesize
\setlength{\tabcolsep}{1.9pt}
\begin{tabular}{llcccccccc}
\toprule
System & Lang & 20 & 15 & 10 & 5 & 2.5 & 0 & $-5$ & $-10$ \\
\midrule
Paraformer-v2 & zh & 0.8 & 0.8 & 0.8 & 2.5 & 9.8 & 42.4 & 101.2 & 101.9 \\
 &  & (0.3,1.3) & (0.3,1.4) & (0.3,1.4) & (1.1,4.2) & (6.8,13.1) & (36.0,49.5) & (98.9,104.1) & (100.2,104.3) \\
 & en & 3.8 & 3.7 & 3.9 & 13.4 & 38.1 & 76.1 & 99.5 & 99.8 \\
 &  & (1.9,6.0) & (1.9,5.9) & (2.1,6.0) & (9.3,17.6) & (32.1,43.8) & (71.3,80.8) & (99.0,99.9) & (99.4,100.0) \\
Qwen-Audio & zh & 0.1 & 0.1 & 0.1 & 0.1 & 0.6 & 8.5 & 127.1 & 138.7 \\
 &  & (0.0,0.2) & (0.0,0.2) & (0.0,0.2) & (0.0,0.2) & (0.1,1.1) & (4.4,13.8) & (110.3,145.0) & (123.6,153.9) \\
 & en & 0.9 & 1.3 & 1.5 & 2.0 & 8.5 & 39.7 & 102.6 & 105.8 \\
 &  & (0.1,2.3) & (0.3,2.7) & (0.4,2.9) & (0.8,3.6) & (3.0,15.1) & (30.3,49.9) & (100.3,105.7) & (100.6,113.3) \\
Qwen3-ASR & zh & 0.1 & 0.2 & 0.1 & 0.3 & 1.0 & 11.2 & 104.8 & 142.9 \\
 &  & (0.0,0.4) & (0.0,0.5) & (0.0,0.4) & (0.0,0.7) & (0.4,1.8) & (7.1,15.6) & (97.2,115.1) & (126.5,163.4) \\
 & en & 0.9 & 0.9 & 0.7 & 2.0 & 4.3 & 34.1 & 101.6 & 104.7 \\
 &  & (0.0,2.2) & (0.1,2.2) & (0.0,1.9) & (0.6,3.8) & (1.7,7.4) & (25.8,43.4) & (98.1,105.5) & (100.0,112.3) \\
\bottomrule
\end{tabular}
\end{table}

\begin{table}[h]
\centering
\caption{Error rate in percent of the streaming translation service
internal frontend recogniser, with session counts; coverage is partial
where the websocket transport dropped sessions.}
\label{tab:ltfrontend}
\footnotesize
\setlength{\tabcolsep}{2.6pt}
\begin{tabular}{lcccccccc}
\toprule
Condition & Clean & 20 dB & 15 dB & 10 dB & 5 dB & 2.5 dB & 0 dB & $-5$ dB \\
\midrule
Error (\%) & 1.0 & 0.8 & 0.2 & 0.9 & 0.0 & 4.0 & 25.4 & 100.8 \\
95\% CI & (0.3,1.7) & (0.3,1.6) & (0.0,0.6) & (0.3,1.6) & (0.0,0.0) & (1.9,6.7) & (19.8,31.4) & (98.4,104.1) \\
$n$ & 144 & 134 & 64 & 144 & 52 & 144 & 144 & 107 \\
\midrule
Condition & Reverb & Narrow & Codec & Loss & Clip & Tempo & $-10$ dB & \\
\midrule
Error (\%) & 0.8 & 0.8 & 0.7 & 0.8 & 0.8 & 1.0 & --- & \\
95\% CI & (0.3,1.5) & (0.3,1.6) & (0.2,1.4) & (0.3,1.5) & (0.3,1.6) & (0.4,1.8) & --- & \\
$n$ & 144 & 144 & 144 & 144 & 142 & 142 & 0 & \\
\bottomrule
\end{tabular}
\end{table}

\section{Response completeness: refusals and empty transcripts}
\label{app:refusal}

Table~\ref{tab:refusal} reports the fraction of utterances for which each
system's final response contains no transcript, split by the two
mechanisms: protocol-level refusals, where the API returns an explicit
no-words error, and silent empties, where a successful response carries an
empty string. The distinction matters operationally: an explicit error can
be caught by the caller and retried or escalated, whereas a silent empty
return is indistinguishable from a genuinely silent audio segment. In the
main grid only the 0~dB condition produces empty outputs; in the sub-zero
extension they become a dominant failure mode for the streaming systems,
with 42 to 58 percent of utterances returning nothing at $-5$~dB, while
the LLM systems increasingly return hallucinated text instead. No failed
API call, of any kind, ever recovered to a different outcome on retry
within the six-attempt budget, and the 24 protocol refusals of the
Qwen-Audio system at 0~dB were stable across repetition runs, so the
behaviour is deterministic in the input rather than an artefact of
transient service conditions.

\begin{table}[h]
\centering
\caption{Empty final transcripts in percent of evaluated utterances, by
mechanism (R: protocol refusal; S: silent empty). Main grid conditions
report the full 480-sentence grid (240 per language); sub-zero conditions
report the 144-sentence sweep subset (72 per language).}
\label{tab:refusal}
\footnotesize
\setlength{\tabcolsep}{2.3pt}
\begin{tabular}{llcccccc}
\toprule
System & Lang & 2.5 dB & 0 dB & $-5$ dB & $-10$ dB & Reverb & Mechanism \\
\midrule
Paraformer-RT-v2 & zh & 4.2 (S) & 2.5 (S) & 41.7 (S) & 43.1 (S) & 0 & silent \\
 & en & 4.2 (S) & 15.0 (S) & 58.3 (S) & 52.8 (S) & 0 & \\
Paraformer-RT-8K & zh & 1.4 (S) & 0.8 (S) & 8.3 (S) & 15.3 (S) & 0.8 (S) & silent \\
Qwen-Audio-3.1 & zh & 0 & 0.8 (R) & 23.6 (R) & 29.2 (R) & 0 & refusal \\
 & en & 0 & 9.2 (R) & 33.3 (R) & 27.8 (R) & 0 & \\
Qwen3-ASR & zh/en & 0 & 0 & 0 & 0 & 0 & none \\
Qwen3.8-Omni & zh/en & 0 & 0 & n/a & n/a & 0 & none \\
\bottomrule
\end{tabular}
\end{table}

\section{Real-speech validation: complete numeric results}
\label{app:real}

Table~\ref{tab:realval} reports the complete recognition results of the
real-speech validation (Section~\ref{sec:results:real}) with bootstrap 95\%
confidence intervals, alongside the empty-output rates at the two hardest
real-babble SNRs. Table~\ref{tab:realladder} reports the six-level
real-babble SNR ladder added in this revision, which densifies the
original three-point ladder ($+5$, $0$, $-5$~dB) with $+10$, $+2.5$, and
$-2.5$~dB points generated with identically seeded babble realisations
from the same forty-source pool, so that the ten-percent crossings
discussed in Section~\ref{sec:results:real} are interpolated between
levels at most 2.5~dB apart. Table~\ref{tab:realval_mt} reports the
corresponding
translation results. The three measured impulse responses used for the
real-reverberation condition are a drama room at MIT campus ($T_{60} =
0.500$~s, DRR $-0.9$~dB), a dining room ($T_{60} = 0.499$~s, DRR
$+1.3$~dB), and a classroom ($T_{60} = 0.495$~s, DRR $+13.1$~dB); for
context, the median reverberation time of the 228 survey responses that
admitted reliable measurement was 0.4~s (interquartile range 0.25--0.62~s)
and the median direct-to-reverberant ratio was $+7.3$~dB (IQR
$+4.1$--$+11.2$~dB).

\begin{table}[t]
\centering
\caption{Real-speech validation: recognition error rate in percent
(CER for Mandarin, WER for English) with bootstrap 95\% confidence
intervals, $n = 75$ per cell. RB denotes real multi-talker babble, TN real
transport noise at 0~dB; Ref.\ gives the empty-output rate in percent at
0~dB/$-5$~dB real babble.}
\label{tab:realval}
\footnotesize
\setlength{\tabcolsep}{2.1pt}
\begin{tabular}{lccccccc}
\toprule
System & Clean & RevR & RB$+5$ & RB$0$ & RB$-5$ & TN$0$ & Ref. \\
\midrule
Qwen-Audio zh & 1.2 (0.2,2.6) & 1.4 (0.5,2.6) & 4.7 (2.0,7.9) & 46.9 (27.5,70.6) & 150.0 (124.4,178.3) & 10.9 (3.8,22.8) & 1/29 \\
\enspace en & 0.8 (0.4,1.3) & 1.6 (0.5,3.2) & 3.1 (1.7,4.8) & 26.0 (19.2,33.5) & 86.6 (82.0,90.8) & 4.6 (2.7,7.3) & 1/19 \\
Qwen3-ASR zh & 1.2 (0.4,2.4) & 1.3 (0.4,2.4) & 15.0 (6.0,27.1) & 117.5 (86.6,150.9) & 207.6 (183.0,235.3) & 8.6 (5.0,13.0) & 0/0 \\
\enspace en & 1.2 (0.5,2.0) & 1.2 (0.6,2.0) & 5.6 (3.7,8.0) & 24.1 (19.0,29.7) & 64.1 (56.9,71.9) & 5.4 (3.5,7.4) & 0/0 \\
Paraformer-v2 zh & 4.1 (2.3,6.3) & 3.2 (1.6,5.1) & 40.1 (32.0,49.1) & 102.6 (88.8,118.0) & 164.4 (145.9,184.2) & 35.6 (29.5,42.4) & 12/7 \\
\enspace en & 3.6 (2.3,5.1) & 4.1 (2.7,5.4) & 24.7 (20.9,28.9) & 56.7 (51.6,61.7) & 86.1 (83.0,89.0) & 30.3 (25.1,35.8) & 0/3 \\
Paraformer-8K zh & 5.1 (3.2,7.5) & 6.3 (4.2,8.8) & 48.6 (42.2,54.8) & 78.7 (74.0,83.4) & 93.0 (89.2,96.8) & 77.4 (71.5,82.6) & 13/27 \\
\bottomrule
\end{tabular}
\end{table}

\begin{table}[t]
\centering
\caption{Real-babble SNR ladder of the real-speech validation, added in
this revision: recognition error rate in percent (CER for Mandarin, WER
for English) with bootstrap 95\% confidence intervals, $n = 75$ per cell,
identically seeded babble realisations from the same forty-source pool as
the original three-point ladder. Values above 100 percent arise from
insertion-dominated outputs.}
\label{tab:realladder}
\footnotesize
\setlength{\tabcolsep}{3.0pt}
\begin{tabular}{llcccccc}
\toprule
System & Lang & $+10$ & $+5$ & $+2.5$ & $0$ & $-2.5$ & $-5$ \\
\midrule
Qwen-Audio & zh & 2.0 & 4.7 & 10.7 & 46.9 & 107.3 & 150.0 \\
 &  & (0.9,3.5) & (2.0,8.1) & (4.3,20.4) & (25.9,71.9) & (81.3,137.9) & (123.7,179.5) \\
 & en & 1.8 & 3.1 & 8.9 & 26.0 & 72.3 & 86.6 \\
 &  & (0.6,3.5) & (1.6,4.9) & (5.4,12.7) & (19.2,33.5) & (63.0,81.0) & (81.8,90.9) \\
Qwen3-ASR & zh & 2.4 & 15.0 & 62.3 & 117.5 & 173.0 & 207.6 \\
 &  & (0.9,4.5) & (5.8,28.7) & (36.2,90.7) & (86.1,152.5) & (138.9,208.1) & (179.0,237.5) \\
 & en & 2.8 & 5.6 & 13.7 & 24.2 & 43.5 & 64.1 \\
 &  & (1.4,4.5) & (3.5,7.9) & (9.5,17.9) & (18.9,29.3) & (35.4,51.1) & (57.0,71.1) \\
Paraformer-v2 & zh & 12.1 & 40.1 & 64.8 & 102.6 & 131.3 & 164.4 \\
 &  & (7.8,17.6) & (31.8,49.1) & (54.9,77.0) & (88.9,118.4) & (116.5,149.2) & (146.3,184.2) \\
 & en & 9.5 & 24.7 & 39.5 & 56.7 & 76.8 & 86.1 \\
 &  & (6.9,12.5) & (20.9,29.1) & (34.4,44.2) & (52.0,61.5) & (72.4,80.9) & (83.0,89.0) \\
Paraformer-8K & zh & 27.9 & 48.6 & 60.2 & 78.7 & 86.8 & 93.1 \\
 &  & (21.9,34.1) & (42.7,54.3) & (54.2,66.5) & (73.8,83.3) & (82.5,90.7) & (89.3,96.8) \\
\bottomrule
\end{tabular}
\end{table}

\begin{table}[t]
\centering
\caption{Real-speech validation: translation quality (corpus BLEU against
the oracle reference produced by the cascade engine on human
ground-truth transcripts; the reference is therefore same-engine with the
cascade and foreign to the end-to-end path, so absolute cross-path gaps
are reference-dependent and are interpreted through the relative-loss
comparison of Section~\ref{sec:results:real}). The end-to-end comparison
is restricted to the three conditions covered by the LiveTranslate quota;
$n = 75$ per cell.}
\label{tab:realval_mt}
\small
\begin{tabular}{lcccccc}
\toprule
 & \multicolumn{2}{c}{Clean} & \multicolumn{2}{c}{Reverb (real)} & \multicolumn{2}{c}{RB $0$~dB} \\
Path & zh$\to$en & en$\to$zh & zh$\to$en & en$\to$zh & zh$\to$en & en$\to$zh \\
\midrule
Cascade (Qwen-Audio $\to$ MT) & 69.2 & 68.9 & 69.7 & 69.8 & 43.7 & 47.8 \\
End-to-end (LiveTranslate) & 44.1 & 39.8 & 44.3 & 38.7 & 20.8 & 28.2 \\
\bottomrule
\end{tabular}
\end{table}

\begin{table}[t]
\centering
\caption{Hold-out recogniser: locally run open-source Whisper small (int8,
CPU) against the commercial systems on the real-speech subset. Error rate
in percent (CER for Mandarin, WER for English) with sentence-level
bootstrap 95\% confidence intervals; ``empty'' counts outputs with no
transcript text. Whisper covers the clean and three babble conditions;
$n = 75$ per language per condition.}
\label{tab:holdout_asr}
\footnotesize
\setlength{\tabcolsep}{2.6pt}
\begin{tabular}{llcccc}
\toprule
System & Lang & Clean & $+5$~dB & $0$~dB & $-5$~dB \\
\midrule
Whisper small & zh & 10.1 & 78.7 & 120.3 & 130.7 \\
 &  & (7.5,13.1) & (60.2,100.8) & (107.5,132.7) & (119.1,144.1) \\
 & en & 3.6 & 19.3 & 47.7 & 92.3 \\
 &  & (2.4,5.1) & (15.3,23.5) & (40.9,55.1) & (86.0,99.1) \\
Qwen-Audio & zh & 1.2 & 4.7 & 46.9 & 150.0 \\
 &  & (0.2,2.5) & (2.0,8.4) & (26.0,71.9) & (124.5,178.5) \\
 & en & 0.8 & 3.1 & 26.0 & 86.6 \\
 &  & (0.4,1.2) & (1.6,4.8) & (19.6,33.3) & (82.3,90.7) \\
Paraformer-v2 & zh & 4.1 & 40.1 & 102.6 & 164.4 \\
 &  & (2.4,6.3) & (32.2,50.2) & (88.3,118.3) & (145.2,184.7) \\
 & en & 3.6 & 24.7 & 56.7 & 86.1 \\
 &  & (2.3,5.0) & (21.0,28.9) & (51.8,62.2) & (82.9,89.2) \\
\bottomrule
\end{tabular}
\end{table}

\begin{table}[t]
\centering
\caption{Measured-RIR reverb calibration: recognition error rate in percent
on the 480-sentence synthetic grid under the original synthetic impulse
response (DRR $-27.6$~dB) and the measured MIT impulse responses (matched
$T_{60} \approx 0.5$~s, DRR $-0.9$/$+1.3$/$+13.1$~dB, rotated). The
cascade BLEU column is scored against the primary reference. The two
end-to-end paths were not re-run (quota exhausted).}
\label{tab:revmit}
\small
\begin{tabular}{lccccc}
\toprule
 & \multicolumn{2}{c}{CER (zh)} & \multicolumn{2}{c}{WER (en)} & BLEU \\
System & synth. & MIT & synth. & MIT & MIT \\
\midrule
Qwen-Audio & 0.2 & 0.2 & 1.2 & 1.1 & 55.0 \\
Qwen3-ASR & 0.1 & 0.2 & 0.5 & 0.5 & 54.9 \\
Paraformer-v2 & 1.3 & 0.7 & 5.1 & 3.7 & 53.3 \\
Paraformer-8K & 18.2 & 3.4 & --- & --- & --- \\
\bottomrule
\end{tabular}
\end{table}

\begin{table}[t]
\centering
\caption{Three-reference scoring of the synthetic-grid cascades: corpus
BLEU (macro-averaged over directions) at clean and 0~dB babble against the
primary (Qwen3.8-Max), oracle (cascade engine), and GLM references, with
the clean-to-0~dB percent relative loss under each. The GLM reference is
family-disjoint from every system and both other references.}
\label{tab:threeref}
\small
\begin{tabular}{lcccccc}
\toprule
 & \multicolumn{3}{c}{Clean BLEU} & \multicolumn{3}{c}{0~dB / rel.\ loss} \\
Cascade & P & O & G & P & O & G \\
\midrule
Qwen-Audio+MT & 53.9 & 89.6 & 51.2 & 44.1 (18\%) & 68.3 (24\%) & 42.4 (17\%) \\
Qwen3-ASR+MT & 53.2 & 94.6 & 52.6 & 43.0 (19\%) & 66.7 (30\%) & 40.6 (23\%) \\
Paraformer+MT & 51.5 & 85.6 & 50.6 & 15.8 (69\%) & 17.4 (80\%) & 13.9 (73\%) \\
\bottomrule
\end{tabular}
\end{table}
